% Source: 04 -Mon Oct 5 2026.pdf, page 13. Content remains in source order.
\sourcepage{04}{13}
The value $\phi(\vec b)$ can be computed in time linear in the formula's size by using the tree.

A satisfying assignment $\vec b$ is such that $\phi(\vec b)=1$. If $\vec b$ exists, $\phi$ is said satisfiable.
\sourceheading{FORMULA SATISFIABILITY (SAT)}
\[
\begin{aligned}
\text{I: }&\langle\phi(x_1,\ldots,x_n)\rangle,\quad\phi\text{ boolean formula},\\
\text{Q: }&\text{is }\phi\text{ satisfiable?}
\end{aligned}
\]
\textbf{EXAMPLE:} $\vec b=(0,0,1,1)$.
\[
\phi(x_1,x_2,x_3,x_4)=(x_1\to x_2)\lor(\neg((\neg x_1\leftrightarrow x_3)\lor x_4)).
\]
\sourcefigure[0.55]{assets/04/page-013-satisfying-formula-tree.png}
$\phi$ is satisfiable $[\phi(\vec b)=1]$.

Very similar to BC-SAT (in fact, this is Cook's original problem).
Nonetheless, the reduction is nontrivial.

1. $\mathrm{SAT}\in\classNP$: Trivial: for $\langle\phi(x_1,\ldots,x_n)\rangle$ a candidate certificate is $\vec b\in\bits^n$. $\phi(\vec b)$ can be computed in linear time using the tree representation.

(write code of VERIFY-SAT$(x,y)$ as an exercise)
